% Part: first-order-logic
% Chapter: natural-deduction
% Section: soundness

\documentclass[../../../include/open-logic-section]{subfiles}

\begin{document}

\iftag{FOL}
      {\olfileid{fol}{ntd}{sou}}
      {\olfileid{pl}{ntd}{sou}}

\olsection{निर्दोषता}

\begin{explain}
नैसर्गिक निगमनासारखी !!^a{derivation} पद्धत \emph{निर्दोष} असते, जर ती
प्रत्यक्षात निष्पन्न न होणाऱ्या गोष्टी !!{derive} करू शकत नसेल. त्यामुळे
निर्दोषता हा !!{derivation} पद्धतींसाठी हमी दिलेल्या सुरक्षिततेचा एक प्रकार
आहे. कोणता सिद्धता-उपपत्तीय गुणधर्म विचारात आहे यावर अवलंबून, उदाहरणार्थ,
आपल्याला पुढील गोष्टी हव्या असतात:
\begin{enumerate}
\item प्रत्येक !!{derivable} !!{sentence} हे \iftag{FOL}{वैध}{सर्वतः सत्य सूत्र} असते;
\item एखादे !!a{sentence} इतर काही वाक्यांपासून !!{derivable} असेल, तर ते
  त्यांचा तार्किक परिणामही असते;
\item !!{sentence}s चा संच विसंगत असेल, तर तो अपूर्ततायोग्य असतो.
\end{enumerate}
हे !!a{derivation} पद्धतीचे महत्त्वाचे गुणधर्म आहेत. यांपैकी एखादा गुणधर्म
खरा नसेल, तर !!{derivation} पद्धतीत दोष आहे---ती खूप जास्त गोष्टी
!!{derive} करेल. म्हणून !!{derivation} पद्धतीची निर्दोषता सिद्ध करणे अत्यंत
महत्त्वाचे आहे.
\end{explain}

\begin{thm}[निर्दोषता]
\ollabel{thm:soundness}
$!A$ हे !!{undischarged} गृहीतके~$\Gamma$ यांच्यापासून !!{derivable} असेल, तर
$\Gamma \Entails !A$.
\end{thm}

\begin{proof}
$\delta$ ही $!A$ ची !!a{derivation} असू द्या. $\delta$ मधील अनुमानांच्या
संख्येवर विगमनाने पुढे जाऊ.

विगमनाच्या आधारपायरीसाठी अनुमानांची संख्या~$0$ असेल तेव्हा दावा सिद्ध करू.
या प्रकरणात $\delta$ मध्ये फक्त एकच !!{sentence}~$!A$, म्हणजे एक गृहीतक,
असते. ते गृहीतक !!{undischarged} असते, कारण गृहीतके फक्त अनुमानांद्वारे
!!{discharged} होऊ शकतात आणि येथे कोणतेही अनुमान नाही. म्हणून सिद्धतेतील
सर्व !!{undischarged} गृहीतके पूर्ण करणारी कोणतीही
\iftag{FOL}{!!{structure}~$\Struct{M}$}{!!{valuation}~$\pAssign{v}$}
$!A$ हेदेखील पूर्ण करते.

आता विगमन पायरी पाहू. $\delta$ मध्ये $n$ अनुमाने आहेत असे समजा. सर्वांत
खालच्या अनुमानाची आधारविधाने अशा उप-!!{derivation}s वापरून !!{derive} केली
आहेत की प्रत्येकामध्ये $n$ पेक्षा कमी अनुमाने आहेत. विगमन गृहीतक असे धरू:
सर्वांत खालच्या अनुमानाची आधारविधाने, त्या आधारविधानांवर संपणाऱ्या
उप-!!{derivation}s च्या !!{undischarged} गृहीतकांपासून निष्पन्न होतात. संपूर्ण
सिद्धतेच्या !!{undischarged} गृहीतकांपासून निष्कर्ष~$!A$ निष्पन्न होतो, हे
आपल्याला दाखवायचे आहे.

सर्वांत खालच्या अनुमानाच्या प्रकारानुसार प्रकरणे वेगळी करू. प्रथम एकच
आधारविधान असलेली संभाव्य अनुमाने पाहू.

\begin{enumerate}
\item शेवटचे अनुमान \Intro{\lnot} आहे असे समजा: !!{derivation} चे रूप
  पुढीलप्रमाणे आहे:
  \begin{prooftree}
    \AxiomC{$\Gamma, \Discharge{!A}{n}$}
    \RightLabel{$\delta_1$}
    \DeduceC{$\lfalse$}
    \DischargeRule{\Intro{\lnot}}{n}
    \UnaryInfC{$\lnot !A$}
  \end{prooftree}
  विगमन गृहीतकानुसार $\lfalse$ हे $\delta_1$ च्या !!{undischarged}
  गृहीतकांपासून~$\Gamma \cup \{!A\}$ निष्पन्न होते. एखादी
  \iftag{FOL}{!!a{structure}~$\Struct{M}$}{!!a{valuation}~$\pAssign{v}$}
  विचारात घ्या. जर $\iftag{FOL}{\Sat{M}}{\pSat{v}}{\Gamma}$, तर
  $\iftag{FOL}{\Sat{M}}{\pSat{v}}{\lnot !A}$, हे दाखवायचे आहे.
  विसंगतीद्वारे सिद्धतेसाठी असे समजा की
  $\iftag{FOL}{\Sat{M}}{\pSat{v}}{\Gamma}$, पण
  $\iftag{FOL}{\Sat/{M}}{\pSat/{v}}{\lnot !A}$, म्हणजे
  $\iftag{FOL}{\Sat{M}}{\pSat{v}}{!A}$. याचा अर्थ
  $\iftag{FOL}{\Sat{M}}{\pSat{v}}{\Gamma \cup \{!A\}}$ असा होईल. हे आपल्या
  विगमन गृहीतकाच्या विरुद्ध आहे. म्हणून
  $\iftag{FOL}{\Sat{M}}{\pSat{v}}{\lnot !A}$.

\item शेवटचे अनुमान \Elim{\land} आहे. याचे दोन प्रकार आहेत: $!A \land
  !B$ या आधारविधानापासून $!A$ किंवा $!B$ निष्पन्न केले जाऊ शकते. पहिले
  प्रकरण पाहू. !!{derivation}~$\delta$ पुढीलप्रमाणे दिसते:
  \begin{prooftree}
    \AxiomC{$\Gamma$}
    \RightLabel{$\delta_1$}
    \DeduceC{$!A \land !B$}
    \RightLabel{\Elim{\land}}
    \UnaryInfC{$!A$}
  \end{prooftree}
  विगमन गृहीतकानुसार $!A \land !B$ हे $\delta_1$ च्या
  !!{undischarged} गृहीतकांपासून~$\Gamma$ निष्पन्न होते. एखादी
  !!a{structure}~\iftag{FOL}{$\Struct{M}$}{$\pAssign{v}$} विचारात घ्या.
  जर $\iftag{FOL}{\Sat{M}}{\pSat{v}}{\Gamma}$, तर
  $\iftag{FOL}{\Sat{M}}{\pSat{v}}{!A}$, हे दाखवायचे आहे.
  $\iftag{FOL}{\Sat{M}}{\pSat{v}}{\Gamma}$ असे समजा. आपल्या विगमन
  गृहीतकानुसार ($\Gamma \Entails !A \land !B$),
  $\iftag{FOL}{\Sat{M}}{\pSat{v}}{!A \land !B}$ हे आपल्याला माहीत आहे.
  व्याख्येनुसार, $\iftag{FOL}{\Sat{M}}{\pSat{v}}{!A \land !B}$ तेव्हा आणि
  केवळ तेव्हाच, जेव्हा $\iftag{FOL}{\Sat{M}}{\pSat{v}}{!A}$ आणि
  $\iftag{FOL}{\Sat{M}}{\pSat{v}}{!B}$. ($!A \land !B$ पासून $!B$
  निष्पन्न करण्याचे प्रकरणही असेच हाताळता येते.)

\item शेवटचे अनुमान \Intro{\lor} आहे. याचे दोन प्रकार आहेत: $!A$ या
  आधारविधानापासून किंवा $!B$ या आधारविधानापासून $!A \lor !B$ निष्पन्न
  केले जाऊ शकते. पहिले प्रकरण पाहू. !!{derivation} चे रूप पुढीलप्रमाणे आहे:
  \begin{prooftree}
    \AxiomC{$\Gamma$}
    \RightLabel{$\delta_1$}
    \DeduceC{$!A$}
    \RightLabel{\Intro{\lor}}
    \UnaryInfC{$!A \lor !B$}
  \end{prooftree}
  विगमन गृहीतकानुसार $!A$ हे $\delta_1$ च्या !!{undischarged}
  गृहीतकांपासून~$\Gamma$ निष्पन्न होते. एखादी
  \iftag{FOL}{!!a{structure}~$\Struct{M}$}{!!a{valuation}~$\pAssign{v}$}
  विचारात घ्या. जर $\iftag{FOL}{\Sat{M}}{\pSat{v}}{\Gamma}$, तर
  $\iftag{FOL}{\Sat{M}}{\pSat{v}}{!A \lor !B}$, हे दाखवायचे आहे.
  $\iftag{FOL}{\Sat{M}}{\pSat{v}}{\Gamma}$ असे समजा; मग
  $\Gamma \Entails !A$ (विगमन गृहीतक) असल्यामुळे
  $\iftag{FOL}{\Sat{M}}{\pSat{v}}{!A}$. म्हणून
  $\iftag{FOL}{\Sat{M}}{\pSat{v}}{!A \lor !B}$ हेदेखील असलेच पाहिजे.
  ($!B$ पासून $!A \lor !B$ निष्पन्न करण्याचे प्रकरणही असेच हाताळता येते.)

\item शेवटचे अनुमान \Intro{\lif} आहे. $!A$ हे गृहीतक आणि $!B$ हा निष्कर्ष
  असलेल्या उपसिद्धतेपासून $!A \lif !B$ निष्पन्न केले जाते, म्हणजे:
  \begin{prooftree}
    \AxiomC{$\Gamma, \Discharge{!A}{n}$}
    \RightLabel{$\delta_1$}
    \DeduceC{$!B$}
    \DischargeRule{\Intro{\lif}}{n}
    \UnaryInfC{$!A \lif !B$}
  \end{prooftree}
  विगमन गृहीतकानुसार $!B$ हे $\delta_1$ च्या !!{undischarged}
  गृहीतकांपासून निष्पन्न होते, म्हणजे $\Gamma \cup \{!A\} \Entails !B$.
  एखादी
  \iftag{FOL}{!!a{structure}~$\Struct{M}$}{!!a{valuation}~$\pAssign{v}$}
  विचारात घ्या. शेवटच्या अनुमानात $!A$ चे विसर्जन झाल्यामुळे $\delta$
  ची !!{undischarged} गृहीतके फक्त $\Gamma$ आहेत. म्हणून
  $\Gamma \Entails !A \lif !B$ हे दाखवायचे आहे. विसंगतीद्वारे सिद्धतेसाठी
  असे समजा की एखाद्या
  \iftag{FOL}{!!{structure}~$\Struct{M}$}{!!{valuation}~$\pAssign{v}$}
  साठी $\iftag{FOL}{\Sat{M}}{\pSat{v}}{\Gamma}$, पण
  $\iftag{FOL}{\Sat/{M}}{\pSat/{v}}{!A \lif !B}$. मग
  $\iftag{FOL}{\Sat{M}}{\pSat{v}}{!A}$ आणि
  $\iftag{FOL}{\Sat/{M}}{\pSat/{v}}{!B}$. पण गृहीतकानुसार $!B$ हा
  $\Gamma \cup \{!A\}$ चा तार्किक परिणाम आहे, म्हणजे
  $\iftag{FOL}{\Sat{M}}{\pSat{v}}{!B}$; ही विसंगती आहे. म्हणून
  $\Gamma \Entails !A \lif !B$.

\item शेवटचे अनुमान \FalseInt आहे. येथे $\delta$ पुढीलप्रमाणे संपते:
  \begin{prooftree}
    \AxiomC{$\Gamma$}
    \RightLabel{$\delta_1$}
    \DeduceC{$\lfalse$}
    \RightLabel{\FalseInt}
    \UnaryInfC{$!A$}
  \end{prooftree}
  विगमन गृहीतकानुसार $\Gamma \Entails \lfalse$. आपल्याला
  $\Gamma \Entails !A$ हे दाखवायचे आहे. तसे नसल्याचे समजा; मग एखाद्या
  $\iftag{FOL}{\Struct{M}}{\pAssign{v}}$ साठी
    $\iftag{FOL}{\Sat{M}}{\pSat{v}}{\Gamma}$ आणि
    $\iftag{FOL}{\Sat/{M}}{\pSat/{v}}{!A}$. पण नेहमीच
    $\iftag{FOL}{\Sat/{M}}{\pSat/{v}}{\lfalse}$ असल्यामुळे याचा अर्थ
    $\Gamma \Entails/ \lfalse$ असा होईल, जो विगमन गृहीतकाच्या विरुद्ध आहे.

\item शेवटचे अनुमान \FalseCl आहे: सराव.

\iftag{FOL}{
\item शेवटचे अनुमान \Intro{\lforall} आहे. मग $\delta$ चे रूप पुढीलप्रमाणे आहे:
  \begin{prooftree}
    \AxiomC{$\Gamma$}
    \RightLabel{$\delta_1$}
    \DeduceC{$!A(a)$}
    \RightLabel{\Intro{\lforall}}
    \UnaryInfC{$\lforall[x][!A(x)]$}
  \end{prooftree}
  विगमन गृहीतकानुसार आधारविधान~$!A(a)$ हे !!{undischarged}
  गृहीतके~$\Gamma$ यांचा तार्किक परिणाम आहे.
  $\Sat{M}{\Gamma}$ अशी एखादी रचना~$\Struct{M}$ विचारात घ्या. आपल्याला
  $\Sat{M}{\lforall[x][!A(x)]}$ हे दाखवायचे आहे.
  $\lforall[x][!A(x)]$ हे !!a{sentence} असल्यामुळे याचा अर्थ प्रत्येक
  चल-विन्यास~$s$ साठी $\Sat{M}{!A(x)}[s]$ दाखवायचे आहे
  (\olref[syn][ass]{prop:sat-quant}). $\Gamma$ मध्ये केवळ वाक्ये असल्यामुळे
  \olref[syn][sat]{defn:satisfaction} नुसार प्रत्येक $!B \in \Gamma$ साठी
  $\Sat{M}{!B}[s]$. $\Struct{M'}$ ही $\Struct{M}$ सारखीच रचना असू द्या,
  फक्त $\Assign{a}{M'} = s(x)$. $a$ हे $\Gamma$ मध्ये येत नसल्यामुळे
  \olref[syn][ext]{cor:extensionality-sent} नुसार
  $\Sat{M'}{\Gamma}$. $\Gamma \Entails !A(a)$ असल्यामुळे
  $\Sat{M'}{!A(a)}$. $!A(a)$ हे !!a{sentence} असल्यामुळे
  \olref[syn][ass]{prop:sentence-sat-true} नुसार
  $\Sat{M'}{!A(a)}[s]$. \olref[syn][ext]{prop:ext-formulas} नुसार
  $\Sat{M'}{!A(x)}[s]$ तेव्हा आणि केवळ तेव्हाच, जेव्हा
  $\Sat{M'}{!A(a)}$ (लक्षात घ्या की $!A(a)$ म्हणजे
  $\Subst{!A(x)}{a}{x}$). म्हणून $\Sat{M'}{!A(x)}[s]$. $a$ हे $!A(x)$
  मध्ये येत नसल्यामुळे \olref[syn][ext]{prop:extensionality} नुसार
  $\Sat{M}{!A(x)}[s]$. पण $s$ हा स्वेच्छ चल-विन्यास होता, म्हणून
  $\Sat{M}{\lforall[x][!A(x)]}$.

\item शेवटचे अनुमान \Intro{\lexists} आहे: सराव.

\item शेवटचे अनुमान \Elim{\forall} आहे: सराव.
}{}
\end{enumerate}

आता अनेक आधारविधाने असलेली संभाव्य अनुमाने पाहू:
\Elim{\lor}, \Intro{\land}, \iftag{FOL}{\Elim{\lif} आणि
  \Elim{\lexists}}{आणि \Elim{\lif}}.
\begin{enumerate}

\item शेवटचे अनुमान \Intro{\land} आहे. $!A$ आणि $!B$ या आधारविधानांपासून
  $!A \land !B$ निष्पन्न केले जाते आणि $\delta$ चे रूप पुढीलप्रमाणे आहे:
  \begin{prooftree}
    \AxiomC{$\Gamma_1$}
    \RightLabel{$\delta_1$}
    \DeduceC{$!A$}
    \AxiomC{$\Gamma_2$}
    \RightLabel{$\delta_2$}
    \DeduceC{$!B$}
    \RightLabel{\Intro{\land}}
    \BinaryInfC{$!A \land !B$}
  \end{prooftree}
  विगमन गृहीतकानुसार $!A$ हे $\delta_1$ च्या !!{undischarged}
  गृहीतकांपासून~$\Gamma_1$ निष्पन्न होते आणि $!B$ हे $\delta_2$ च्या
  !!{undischarged} गृहीतकांपासून~$\Gamma_2$ निष्पन्न होते. $\delta$ ची
  !!{undischarged} गृहीतके $\Gamma_1 \cup \Gamma_2$ आहेत; म्हणून
  $\Gamma_1 \cup \Gamma_2 \Entails !A \land !B$ हे दाखवायचे आहे.
  $\iftag{FOL}{\Sat{M}}{\pSat{v}}{\Gamma_1 \cup \Gamma_2}$ असलेली एखादी
  \iftag{FOL}{!!a{structure}~$\Struct{M}$}{!!a{valuation}~$\pAssign{v}$}
  विचारात घ्या. $\iftag{FOL}{\Sat{M}}{\pSat{v}}{\Gamma_1}$ आणि
  $\Gamma_1 \Entails !A$ असल्यामुळे
  $\iftag{FOL}{\Sat{M}}{\pSat{v}}{!A}$ असलेच पाहिजे; तसेच
  $\iftag{FOL}{\Sat{M}}{\pSat{v}}{\Gamma_2}$ आणि
  $\Gamma_2 \Entails !B$ असल्यामुळे
  $\iftag{FOL}{\Sat{M}}{\pSat{v}}{!B}$. या दोन्हींवरून
  $\iftag{FOL}{\Sat{M}}{\pSat{v}}{!A \land !B}$.

\item शेवटचे अनुमान \Elim{\lor} आहे: सराव.

\item शेवटचे अनुमान \Elim{\lif} आहे. $!A \lif !B$ आणि $!A$ या
  आधारविधानांपासून $!B$ निष्पन्न केले जाते. !!{derivation}~$\delta$
  पुढीलप्रमाणे दिसते:
  \begin{prooftree}
    \AxiomC{$\Gamma_1$}
    \RightLabel{$\delta_1$}
    \DeduceC{$!A \lif !B$}
    \AxiomC{$\Gamma_2$}
    \RightLabel{$\delta_2$}
    \DeduceC{$!A$}
    \RightLabel{\Elim{\lif}}
    \BinaryInfC{$!B$}
  \end{prooftree}
  विगमन गृहीतकानुसार $!A \lif !B$ हे $\delta_1$ च्या
  !!{undischarged} गृहीतकांपासून~$\Gamma_1$ निष्पन्न होते आणि $!A$ हे
  $\delta_2$ च्या !!{undischarged} गृहीतकांपासून~$\Gamma_2$ निष्पन्न होते.
  एखादी
  \iftag{FOL}{!!a{structure}~$\Struct{M}$}{!!a{valuation}~$\pAssign{v}$}
  विचारात घ्या. जर $\iftag{FOL}{\Sat{M}}{\pSat{v}}{\Gamma_1 \cup
    \Gamma_2}$, तर $\iftag{FOL}{\Sat{M}}{\pSat{v}}{!B}$, हे दाखवायचे आहे.
  $\iftag{FOL}{\Sat{M}}{\pSat{v}}{\Gamma_1 \cup \Gamma_2}$ असे समजा.
  $\Gamma_1 \Entails !A \lif !B$ असल्यामुळे
  $\iftag{FOL}{\Sat{M}}{\pSat{v}}{!A \lif !B}$. तसेच
  $\Gamma_2 \Entails !A$ असल्यामुळे
  $\iftag{FOL}{\Sat{M}}{\pSat{v}}{!A}$. याचा अर्थ
  $\iftag{FOL}{\Sat{M}}{\pSat{v}}{!B}$ असा होतो (कारण
  $\iftag{FOL}{\Sat/{M}}{\pSat/{v}}{!B}$ असेल, तर
  $\iftag{FOL}{\Sat{M}}{\pSat{v}}{!A}$ असल्यामुळे
  $\iftag{FOL}{\Sat/{M}}{\pSat/{v}}{!A \lif !B}$ असे होईल, जे
  $\iftag{FOL}{\Sat{M}}{\pSat{v}}{!A \lif !B}$ च्या विरुद्ध आहे).

\item शेवटचे अनुमान \Elim{\lnot} आहे: सराव.

\tagitem{FOL}{शेवटचे अनुमान \Elim{\lexists} आहे: सराव.}{}
\end{enumerate}
\end{proof}

\tagprob{FOL}
\begin{prob}
\olref[fol][ntd][sou]{thm:soundness} ची सिद्धता पूर्ण करा.
\end{prob}
\tagendprob

\tagprob{notFOL}
\begin{prob}
\olref[pl][ntd][sou]{thm:soundness} ची सिद्धता पूर्ण करा.
\end{prob}
\tagendprob

\begin{cor}
\ollabel{cor:weak-soundness}
जर $\Proves !A$, तर $!A$ हे \iftag{FOL}{वैध}{सर्वतः सत्य सूत्र} आहे.
\end{cor}

\begin{cor}
\ollabel{cor:consistency-soundness}
जर $\Gamma$ पूर्ततायोग्य असेल, तर तो सुसंगत आहे.
\end{cor}

\begin{proof}
प्रतिलोम विधान सिद्ध करू. $\Gamma$ सुसंगत नाही असे समजा. मग
$\Gamma \Proves \lfalse$, म्हणजे $\Gamma$ मधील !!{undischarged}
गृहीतकांपासून $\lfalse$ ची !!a{derivation} आहे. \olref{thm:soundness}
नुसार $\Gamma$ पूर्ण करणारी कोणतीही
\iftag{FOL}{!!{structure}~$\Struct{M}$}{!!{valuation}~$\pAssign{v}$}
$\lfalse$ देखील पूर्ण करते. प्रत्येक
\iftag{FOL}{!!{structure}~$\Struct{M}$}{!!{valuation}~$\pAssign{v}$}
साठी $\iftag{FOL}{\Sat/{M}}{\pSat/{v}}{\lfalse}$ असल्यामुळे कोणतीही
\iftag{FOL}{$\Struct{M}$}{$\pAssign{v}$} $\Gamma$ पूर्ण करू शकत नाही,
म्हणजे $\Gamma$ पूर्ततायोग्य नाही.
\end{proof}

\end{document}
